% PF22091C
% Wahnvorstellungen
% 14.0
% 2013-11-30
:ref:`m-wahnvorstellungen`
Die -- für
den Außenstehenden -- wirren Vorstellungen stellen das Fachpersonal vor die Frage, wie mit diesen
Patienten umzugehen ist.
Ziel ist:
\begin{enumerate}
  -   Abwendung von unmittelbaren Gefahren für den
  Patienten und die Umgebung
  -   Korrekte und professionelle Durchführung eines
  Transportes in eine geeignete Einrichtung.
\end{enumerate}
Es gelten somit folgende Grundsätze:
\begin{enumerate}
  -   Sofern eine Gefährdung für den Patienten oder für
  Dritte besteht, ist diese (wenn möglich) zu beseitigen,
  dazu zählt u. U. auch der Rückzug und die Herbeiziehung
  der Exekutive.
  -   Wahnvorstellungen sollen **nicht ausgeredet**
  werden.
  -    Wahnvorstellungen sollen aber auch **nicht
  bestätigt** werden.
  -   Am besten man nimmt die Schilderungen des
  Patienten **interessiert, aber neutral und möglichst
  kommentarlos** zur Kenntnis (\emph{\enquote{Mhm!}, \enquote{Aha.}}).
\end{enumerate}